\documentclass[a4paper, 11pt]{article}

% fontspec allows you to use TTF/OTF fonts directly (XeLaTeX/LuaLaTeX only)
% \usepackage{fontspec}
% \defaultfontfeatures{Ligatures=TeX}
% \setmainfont[Path=fonts/, BoldFont=Outfit-Medium]{Outfit-Light}
\renewcommand{\familydefault}{\sfdefault}

% Icons
\usepackage{fontawesome5}
\usepackage{academicons}

% Colours and graphics
\usepackage{graphicx}
\usepackage{xcolor}
\definecolor{bright_blue}{HTML}{0CC6DE}
\definecolor{dark_blue}{HTML}{002F5F}
\definecolor{grey}{HTML}{575757}
\newcommand{\highlightbold}[1]{\textbf{\textcolor{bright_blue}{#1}}}
\newcommand{\highlightdark}[1]{\textcolor{dark_blue}{#1}}
\newcommand{\tbc}[1]{\textcolor{grey}{#1}}

% Geometry
\usepackage{geometry}
\geometry{a4paper, top=10mm, bottom=10mm, left=10mm, right=10mm, heightrounded, includehead, headheight=18.5pt}

% Headers
\usepackage{lastpage}
\usepackage{fancyhdr}
\pagestyle{fancy}
\fancyhead{}
\fancyfoot{}
\fancyhead[C]{\scriptsize\textcolor{dark_blue}{Denis Sergeev's CV\\\thepage~/~\pageref*{LastPage}}}
\renewcommand{\headrulewidth}{0pt}

% Paragraph formatting
\usepackage{parskip}

% Tables
\usepackage{tabularx}
\usepackage{ltablex}
\usepackage{multirow}

% Section formatting
\usepackage{titlesec}
\titleformat{\section}{\Large\bfseries\raggedright}{}{0em}{}[\titlerule]
\titlespacing{\section}{0pt}{5pt}{0pt}
\titlespacing*{\subsection}{0pt}{5pt}{0pt}

% Links
\usepackage[unicode, draft=false]{hyperref}
\hypersetup{
colorlinks,
breaklinks,
urlcolor=dark_blue,
linkcolor=dark_blue,
pdfauthor={Denis Sergeev},
pdftitle={Denis Sergeev's CV}
}

% Date formatting
\usepackage[nodayofweek]{datetime}
\newdateformat{mydate}{\THEDAY\ \shortmonthname[\THEMONTH]\ \THEYEAR}

% Units formatting
\usepackage{siunitx}
\sisetup{group-separator={\,}}
\newcommand{\pval}[1]{\pounds\num{#1}}
\newcommand{\estval}[1]{$\sim$\pounds\num{#1}}

% Underlining
\usepackage{soul}

\newcommand{\sep}{~~$\circ$~~}

\newcommand{\link}[1]{{\href{#1}{\scriptsize\faExternalLink*}}}

\newcommand{\no}{{\scriptsize\#}}

% To debug page outer frames
% \usepackage{showframe}



\begin{document}
\thispagestyle{empty}

% Needed to keep the tabularx columns stretched over the whole linewidth when using ltablex
\keepXColumns

\noindent\vspace{-20mm}
\begin{tabularx}{\linewidth}{@{}p{0.25\linewidth} p{0.55\linewidth} X@{}}
\multicolumn{3}{@{}l@{}}{{\Huge\textbf{\textcolor{exeter_bright_green}{Denis Sergeev}}}} \\
& & \\
\href{https://pronouns.org/what-and-why}{\,\faUser~Pronouns: he/him/his} &
\multirow{5}{*}{%
\begin{minipage}[t]{\linewidth}
\includegraphics[width=\linewidth]{images/sergeev_ads_metrics.pdf}
\end{minipage}
} &
\multirow{5}{*}{%
\begin{minipage}[t]{\linewidth}
\begin{tabular}{@{}r l@{}}
\input{stats}
 & \\
\multicolumn{2}{r}{\tiny \href{https://ui.adsabs.harvard.edu/public-libraries/2IT653szTA-gYw3vyA-9Sw}{\aiADS\,Updated:\,\mydate\today}}
\end{tabular}
\end{minipage}
} \\
\href{https://www.bristol.ac.uk/physics/research/astrophysics}{\faHome~University of Bristol, UK} & &\\
\href{mailto:denis.sergeev@bristol.ac.uk}{\faEnvelope~denis.sergeev@bristol.ac.uk} & &\\
\href{http://orcid.org/0000-0001-8832-5288}{\aiOrcid~0000-0001-8832-5288} & &\\
\href{https://dennissergeev.github.io}{\faGlobe~dennissergeev.github.io} & &\\
\href{https://github.com/dennissergeev}{\faGithub~dennissergeev} & &\\
\end{tabularx}
\vspace{-5mm}


% \section{About}
% I am fascinated by atmospheric phenomena on different planets.
% I am especially interested in the role of convection and clouds in planetary climates and how cloud processes affect the habitability of terrestrial exoplanets.
% In my research, I try to unravel the complexity of planetary atmospheres by using 3D climate models.
% \section{Research Interests}
% \vspace{10pt}
% Planetary Climates\sep Exoplanets\sep Atmospheric Convection\sep Cloud Microphysics\sep General Circulation Models


\section{Career history}
\begin{tabularx}{\linewidth}{@{}r X@{}}
Jan 2025--now  & \textbf{Lecturer in Astrophysics}\\
& School of Physics, University of Bristol \\
Sep 2021--Dec 2024 & \textbf{Postdoctoral Researcher}\\
& Project: Exascale Exoplanet Modelling\\
& Department of Physics \& Astronomy, University of Exeter \\
Sep 2018--Aug 2021 & \textbf{Postdoctoral Researcher} \\
& Project: Climate Modelling of Rocky Exoplanets\\
& Department of Mathematics \& Statistics, University of Exeter
\end{tabularx}

\vspace{-10pt}

\section{Academic Qualifications}
\begin{tabularx}{\linewidth}{@{}r X@{}}
Oct 2014--Aug 2018 & \textbf{PhD in Meteorology}\\
& Thesis title (shortened): `Characteristics of Polar Lows in the Nordic Seas'~\link{https://ueaeprints.uea.ac.uk/id/eprint/68204/} \\ %  and the Impact of Orography and Sea Ice on Their Development
& School of Environmental Sciences, University of East Anglia \\
& Supervisors: Ian A. Renfrew, Thomas Spengler, Stephen Dorling \\
Sep 2009--May 2014 & \textbf{Specialist Diploma} \\
& Thesis title: `Idealised Numerical Modelling of Polar Mesocyclone Dynamics'~\link{https://figshare.com/articles/thesis/SergeevDE_diploma_pdf/5326846/1}\\
& Department of Meteorology and Climatology, Moscow State University  \\
& Supervisor: Victor Stepanenko
\end{tabularx}


\section{Funding and Awards}
\begin{tabularx}{\linewidth}{@{}lXr@{}}
\multicolumn{2}{@{}l}{\highlightbold{Direct Funding, PI}} & \highlightbold{Est. Total Value} \\
2026 & Undergraduate Student Bursary \sep RAS & \pval{1500} \\
2026 & Summer Student Placement Bursary \sep University of Bristol & \estval{4500} \\
2024 & Above \& Beyond Silver Award \sep University of Exeter & \pval{1000} \\
2023 & Meeting Organisation Funding (Exoclimes VI and ExoSLAM) \sep RAS & \pval{5000} \\
2022 & Undergraduate Student Bursary (awarded; student declined) \sep RAS & \pval{1200} \\
2017 & Best Presentation Award \sep CEEDA Symposium & \estval{100} \\
2016 & Travel Bursary \sep Polar Prediction School & \estval{1000} \\
2015 & Travel Award \sep High-Latitude Dynamics workshop & \estval{1000} \\ % https://highlatdynamics.w.uib.no/
2014 & Lord Zuckerman PhD scholarship \sep School of Environmental Sciences, UEA & \estval{112000} \\
2014 & Young Scientist Travel Award \sep EGU General Assembly & \estval{200} \\
2014 & Russian Academy of Sciences Young Scientist Medal & \estval{1000} \\
\multicolumn{2}{@{}l}{\highlightbold{Direct Funding, co-I}} & \\
2025 & Meeting Organisation Funding (UKExoM 2026) \sep RAS & \pval{2500} \\
2025 & Isambard 3 Allocation (30,000 node-hours; PI: N.J. Mayne) \sep UKRI (Isambard 3) & \textbullet \\
2024 & UKSA Studentships: Mars Exploration Science & \estval{100000} \\
2024 & Research Software Engineer Support \sep DiRAC HPC & \estval{45000} \\
\multicolumn{2}{@{}l}{\highlightbold{Observational Facilities Resources, co-I}} & \\
2023 & JWST Cycle 2, 49.21 hours (GO 3838~\link{https://www.stsci.edu/jwst/science-execution/program-information?id=3838}, PI: J. Kirk) & \textbullet \\
\end{tabularx}


% \section{Research Interests}
% \vspace{-5pt}
% \begin{tabularx}{\linewidth}{@{}Xl@{}}
% & {\scriptsize Publications (see below):}\\
% \noindent\highlightbold{Atmospheric aerosols:} & \no17, 20, 22, 24, 27, 29, 34 \\
% \sep How do clouds, hazes and dust shape planetary climates? & \\
% %
% \noindent\highlightbold{Atmospheric convection on exoplanets:} & \no8, 17, 30, 35 \\
% \sep How does convection shape global energy redistribution? & \\
% %
% \noindent\highlightbold{Extraterrestrial lightning:} & \no19, 35 \\
% \sep How is lightning generated on exoplanets and can we detect it? & \\
% %
% \noindent\highlightbold{Atmospheric dynamics on Earth and other planets:} & \no2, 3, 4, 5, 7, 10, 11, 13, 14, 16 \\
% \sep How do wind jets and cyclones form in planetary atmospheres? & \\
% %
% \noindent\highlightbold{Planet formation and atmospheric evolution:}
% & \no25, 28, 31, 32, 33, 36, 37, 38 \\
% \sep How does atmospheric composition relate to the planet's history? & \\
% %
% \noindent\highlightbold{Model development and intercomparison:} & \no12, 15, 16, 17, 18, 23, 30, 36 \\
% \sep How do we build robust and reproducible exoplanet models? & \\
% %
% \end{tabularx}


\section{Publications}
\begin{tabularx}{\linewidth}{@{}rXl@{}}
\no & {\scriptsize\tbc{(preprints in \textbf{grey})}} & {\scriptsize\highlightdark{Citations}}\\
\input{pubs}
\end{tabularx}


\section{Conferences and Seminars}
\vspace{10pt}
\subsection*{Invited Talks}
\begin{tabularx}{\linewidth}{@{}l X@{}}
May 2025 & General circulation models for exoplanet atmospheres \\ & JpGU-AGU Joint Meeting 2026 \sep Chiba, Japan \\
Oct 2025 & CUISINES - a framework for exoplanet model intercomparison projects \\ & Atmospheric and interior evolution of planetary magma oceans \sep Leiden, the Netherlands \\
Jun 2025 & Atmospheric dynamics on other planets~\link{https://durham.readthedocs.io/en/latest/hpcdays/workshops.html} \\ & Durham HPC Days \sep Durham, UK \\
Feb 2025 & Exoplanet climate modelling with LFRic \\ & University of East Anglia \sep Norwich, UK \\
May 2024 & 3D simulations of exoplanet atmospheres with the next-generation Met Office model \\ & University of Leicester \sep Leicester, UK \\
Apr 2024 & Shall I compare thee to a distant world? Inter-planet and inter-model comparative studies \\ & EGU General Assembly \sep Vienna, Austria \\
Jul 2023 & Simulations of idealised 3D atmospheric flows on terrestrial planets using LFRic-Atmosphere \\ & NASA GISS Seminar \sep Online \\
Mar 2023 & First results of using LFRic for exoplanet climate modelling \\ & NIWA Seminar \sep Wellington, New Zealand \\
Feb 2023 & Atmospheric dynamics and chemistry on exoplanets \\ & UQ Astro Group Meeting \sep Brisbane, Australia \\
Feb 2023 & Atmospheric dynamics and chemistry on exoplanets~\link{https://youtu.be/g4d-ibZvBIY} \\ & UniSQ Exoplanet Group Seminar \sep Brisbane, Australia \\
Feb 2023 & Atmospheric dynamics and chemistry on exoplanets \\ & UNSW AstroSeminar \sep Sydney, Australia \\
Apr 2022 & Dichotomy of the atmospheric circulation on TRAPPIST-1e~\link{https://youtu.be/0uDBIp_EQrg} \\ & NASA GISS Seminar \sep Online \\
Jan 2022 & Dichotomy of the atmospheric circulation on TRAPPIST-1e \\ & NASA GSFC Extrasolar Planets Seminar \sep Online \\
Nov 2021 & TRAPPIST-1 Habitable Atmosphere Intercomparison (THAI) \\ & MPIA APEx Exocoffee \sep Online \\
May 2021 & Overcast on TRAPPIST-1e~\link{https://youtu.be/ZUfIK-HMgLw} \\ & RCC MSU Geophysical Seminar \sep Online \\
Sep 2020 & Simulations of convection over a range of atmospheric conditions on TRAPPIST-1e~\link{https://youtu.be/E4UAoCoI1x8} \\ & THAI Workshop \sep Online \\
Apr 2020 & Atmospheric convection plays a key role in the climate of tidally locked exoplanets~\link{https://slides.com/denissergeev/2020-04-27-uor-met} \\ & University of Reading Meteorology Seminar \sep Online \\
Apr 2020 & Atmospheric convection plays a key role in the climate of tidally locked exoplanets~\link{https://slides.com/denissergeev/2020-04-21-nasa-giss} \\ & NASA GISS Seminar \sep Online \\
\end{tabularx}
\subsection*{Contributed Talks}
\begin{tabularx}{\linewidth}{@{}l X@{}}
May 2026 & Lightning climatology on TRAPPIST-1e \\ & ELSI workshop \sep Tokyo, Japan \\
Sep 2023 & Introducing GeoVista - Cartographic rendering and mesh analytics powered by PyVista (joint talk) \\ & Met Office Seminar \sep Exeter, UK \\
Jul 2022 & Bistability of the atmospheric circulation on TRAPPIST-1e \\ & Rocky Worlds II \sep Oxford, UK \\
Apr 2022 & Dichotomy of the atmospheric circulation on TRAPPIST-1e \\ & Exoplanet Modelling in the James Webb Era II: Terrestrial planets and sub-Neptunes \sep Online \\
Nov 2020 & Explicit convection on tidally locked rocky exoplanets simulated with the UM nesting suite~\link{https://slides.com/denissergeev/2020-11-10-um-workshop-sergeev} \\ & Unified Model users workshop \sep Online \\
Aug 2019 & Simulations of moist convection on tidally-locked rocky exoplanets~\link{https://youtu.be/9nGIpQiPwDs} \\ & Exoclimes V \sep Oxford, UK \\
Jun 2019 & North Atlantic polar mesoscale cyclones in ERA5 and ERA-Interim reanalyses~\link{https://speakerdeck.com/dennissergeev/north-atlantic-polar-mesoscale-cyclones-in-era5-and-era-interim-reanalyses} \\ & IGP workshop \sep Norwich, UK \\
Apr 2019 & Atmospheric convection on tidally-locked Earth-like exoplanets \\ & UK Exoplanet Community Meeting \sep London, UK \\
Jun 2018 & Modification of Polar Low Development by Sea Ice and Svalbard Orography~\link{https://speakerdeck.com/dennissergeev/polar2018} \\ & POLAR2018 \sep Davos, Switzerland \\
Oct 2017 & The influence of Svalbard orography and sea ice on polar low development~\link{https://figshare.com/articles/The_influence_of_Svalbard_orography_and_sea_ice_on_polar_low_development/5510416} \\ & 18th Cyclone Workshop \sep Sainte-Adèle, Canada \\
Apr 2017 & Polar lows and how background environment can influence their development~\link{http://dennissergeev.github.io/ceeda2017} \\ & Cambridge Earth Systems Science EnvEast Doctoral Alliance Symposium \sep Cambridge, UK \\
May 2016 & Structure of the shear-line polar low south of Svalbard \\ & NORPAN meeting \sep Tokyo, Japan \\
Apr 2016 & Structure of the shear-line polar low south of Svalbard~\link{http://slides.com/denissergeev/deck} \\ & 13th European Polar Lows Working Group Workshop \sep Paris, France \\
\end{tabularx}
\subsection*{Poster Presentations}
\begin{tabularx}{\linewidth}{@{}l X@{}}
Nov 2025 & Lightning climatology on rocky exoplanets in a global storm-resolving model \\ & CTR Wilson Meeting on Atmospheric Electricity \sep Bath, UK \\
Jun 2024 & The impact of convection on the climate of a tidally locked planet in stretched-mesh simulations \\ & Exoplanets 5 \sep Leiden, Netherlands \\
Apr 2024 & The impact of convection on the climate of TRAPPIST-1e in global stretched-mesh simulations \\ & EGU General Assembly \sep Vienna, Austria \\
Apr 2024 & The impact of convection on the climate of a tidally locked planet in stretched-mesh simulations \\ & UK Exoplanet Community Meeting \sep Birmingham, UK \\
Nov 2022 & Dry Modern-Day Mars Climate in the Met Office Unified Model \\ & UK Solar System Planetary Atmospheres \sep London, UK \\
Sep 2022 & Bistability of the Atmospheric Circulation on TRAPPIST-1e \\ & UK Exoplanet Community Meeting \sep Edinburgh, UK \\
Jul 2015 & Structure and dynamics of a shear-line polar low during a cold-air outbreak over the Norwegian Sea \\ & Royal Meteorological Society Student Conference \sep Birmingham, UK \\
Mar 2015 & Structure and dynamics of a shear-line polar low during a cold-air outbreak over the Norwegian Sea \\ & Dynamics of Atmosphere-Ice-Ocean Interactions in the High Latitudes workshop \sep Rosendal, Norway \\
May 2014 & Numerical modelling of polar mesocyclones dynamics diagnosed by the energy budget \\ & EGU General Assembly \sep Vienna, Austria \\
Apr 2013 & Impact of subgrid-scale vegetation heterogeneity on the carbon cycle \\ & EGU General Assembly \sep Vienna, Austria \\
Apr 2013 & Numerical modelling of polar mesocyclones generation mechanisms \\ & EGU General Assembly \sep Vienna, Austria \\
\end{tabularx}


\section{Supervision}
\vspace{5pt}
{\scriptsize(Projects with me as the lead supervisor are in \highlightbold{bold}. Students who continued their academic career are \ul{underlined}.)}
\subsection*{PhD Supervision}
\begin{tabularx}{\linewidth}{@{}l X@{}}
Sep 2025--Sep 2029 & \highlightbold{Alex Corbett} \\
                   & Project: \highlightbold{Convection on Sub-Neptunes} \\
                   & Co-supervisors: B. Shipway, Z. Leinhardt \\
Sep 2025--Sep 2029 & Will Luscombe (U. Exeter) \\
                   & Project: Forecasting Martian dust storms \\
                   & Co-supervisors: N. J. Mayne, M. Bate, B. Drummond \\
Sep 2021--Apr 2025 & \ul{Mei Ting (Martha) Mak} (U. Exeter) \\
                   & Project: Hazes in Planetary Atmospheres \\
                   & Co-supervisors: N. J. Mayne, J. Manners, E. Hébrard \\
\end{tabularx}
\vspace{-10pt}
\subsection*{MSc and MRes Supervision}
\begin{tabularx}{\linewidth}{@{}l X@{}}
Jan 2027--Dec 2029 & \highlightbold{Grace Casson} (MSci by Research) \\
                   & Project: \highlightbold{Distinguishing Global Cloud Regimes on Earth-like Exoplanets} \\
                   & Co-supervisors: V. Kofman (U. Oslo) \\
Sep 2026--May 2027 & \highlightbold{Cleo Hyde \& Lily Plummer} \\
                   & Project: \highlightbold{Ice Giant Atmospheres in 3D} \\
Sep 2026--May 2027 & \highlightbold{Matthew Pickering} \\
                   & Project: \highlightbold{Modelling Convection in Hydrogen-Rich Atmospheres} \\
                   & Co-supervisors: A. Corbett\\
Sep 2025--May 2026 & \highlightbold{\ul{Freya Evans} \& Daisy Green} \\
                   & Project: \highlightbold{Atmospheric Dynamics on Ice Giants} \\
Sep 2025--May 2026 & \highlightbold{Catherine Kerr \& Lily Odhuba} \\
                   & Project: \highlightbold{Lightning Storms on Earth-like Exoplanets} \\
Jan 2023--May 2024 & Tom Batchelor, Luke Benzing, \& \ul{Alex McGinty} \\
                   & Project: Mars Atmosphere Modelling \\
                   & Co-supervisors: M. Bate, N. J. Mayne, D. McCulloch \\
Sep 2020--Sep 2022 & \ul{Danny McCulloch} (MSci by Research) \\
                   & Project: Climate Modelling of Modern-Day Mars \\
                   & Co-supervisors: M. Bate, N. J. Mayne \\
Apr 2021--Sep 2022 & \ul{Meghan Plumridge} (MSci by Research) \\
                   & Project: Climate Modelling of Early Mars \\
                   & Co-supervisors: M. Bate, N. J. Mayne \\
Jan 2021--May 2022 & Jasper Chadwick \& Esse Sellwood \\
                   & Project: Ocean Heat Transport on Rocky Exoplanets \\
                   & Co-supervisors: F. H. Lambert, J. Eager-Nash \\
Jan 2021--May 2022 & Isabelle Browne \& \ul{Oakley Young} \\
                   & Project: Greenhouse Effect on Early Mars \\
                   & Co-supervisors: F. H. Lambert, N. J. Mayne, J. Eager-Nash \\
Jan 2020--May 2021 & Toby Ferrison \\
                   & Project: Titan Climate Modelling \\
                   & Co-supervisor: F. H. Lambert \\
Oct 2018--May 2019 & \ul{Jake Eager-Nash} \& David Reichelt \\
                   & Project: Implications of Stellar Type on the Climate of Tidally Locked Terrestrial Exoplanets \\
                   & Co-supervisors: F. H. Lambert, N. J. Mayne \\
\end{tabularx}
\vspace{-10pt}
\subsection*{Undergraduate and Summer Internship Supervision}
\begin{tabularx}{\linewidth}{@{}l X@{}}
Jun--Aug 2026 & \highlightbold{Evie Cushing} \\
              & Project: \highlightbold{Radiative Transfer Modelling for Ice Giants} \\
Jul--Sep 2022 & \ul{Oakley Young} \\
              & Project: Ekman Ocean Model \\
              & Co-supervisors: J. Eager-Nash, F. H. Lambert \\
Jun--Sep 2022 & \highlightbold{\ul{James McDermott} \& \ul{Lottie Woods}} \\
              & Project: \highlightbold{Simulations of Lightning Storms on Tidally Locked Rocky Exoplanets} \\
Jun--Aug 2021 & \ul{Oakley Young} \\
              & Project: Climate Modelling of Archean Earth \\
              & Co-supervisors: J. Eager-Nash, N. J. Mayne \\
Jun--Aug 2021 & Joshua Parkin \& Esse Sellwood \\
              & Project: The Impact of Host Star Spectrum on the Climate of Rocky Exoplanets \\
              & Co-supervisors: J. Eager-Nash, N. J. Mayne \\
Jun--Aug 2019 & \ul{Isobel Parry} \\
              & Project: Water Cycle on Proxima Centauri b \\
              & Co-supervisor: F. H. Lambert \\
\end{tabularx}

\section{PhD Examinations}
\begin{tabularx}{\linewidth}{@{}l X@{}}
2025 & Maria Di Paolo \sep University of East Anglia \sep Supervisors: Manoj Joshi, David Stevens \\
\end{tabularx}

\section{Teaching and Mentoring}
\begin{tabularx}{\linewidth}{@{}r X@{}}
2026--now & Fellow of the Higher Education Academy (AdvanceHE UK) \\
2026--now & Environmental Physics \\ & Lecturer \sep University of Bristol \\
2025--now & Practical Physics III: Research Skills and Group Project \\ & Tutor \sep University of Bristol \\
2025--now & Research Project in Physics \\ & Supervisor \& assessor \sep University of Bristol \\
Jul 2024 & Algorithms For Exascale Summer School~\link{https://sites.google.com/view/algorithmsforexascale/home} \\ & Invited lecturer \sep University of Exeter \\
Feb 2024 & Physics of Climate Change \\ & Workshop lead \sep University of Exeter\\
Jul 2023 & Climatematch Academy \\ & Mentor \sep Online\\
Jul 2023 & International Sustainability Summer School \\ & Lecturer \sep University of Exeter \\
Jun 2023 & Exoclimes Summer School in Atmospheres and Modelling (ExoSLAM)~\link{https://exoclimes.org/exoslam/index.html} \\ & Lecturer \sep University of Exeter \\
2016--2018 & Introduction to Python in Environmental Sciences~\link{https://ueapy.github.io/pythoncourse2018} \\ & Course creator \& lead \sep University of East Anglia \\
2015--2017 & Modelling Environmental Processes; Meteorology; Numerical Skills \\ & Teaching assistant \sep University of East Anglia \\
\end{tabularx}


\section{Active large collaborations}
\begin{tabularx}{\linewidth}{@{}l X@{}}
2024--\hphantom{99} & \textbf{Co-lead}\sep CUISINES, an international network for exoplanet model intercomparison projects~\link{https://nexss.info/cuisines}\\
2022-- & \textbf{Member}\sep BOWIE+, a UK-based network for studying exoplanet atmospheres with JWST~\link{https://bowiecollaboration.github.io} \\
2018-- & \textbf{Member, co-lead from 2025}\sep Exoclimatology Theory Group~\link{https://exoclimatology.com} 
\end{tabularx}

\section{Organisation of scientific meetings}
\begin{tabularx}{\linewidth}{@{}l X@{}}
% 2024--\hphantom{99} & \ul{Co-lead}\sep CUISINES~\link{https://nexss.info/cuisines}, an international network for exoplanet model intercomparison projects\\
% 2024-- & \ul{Main organiser}\sep BUFFET meetings \\
2027\hphantom{--99} & \textbf{Main organiser}\sep Bridging atmospheres and interiors: Combined modelling and experimental approaches for sub-Neptune exoplanets (RAS Specialist Discussion)\sep London, UK \\
2026 & \textbf{Main organiser}\sep BUFFET-6~\link{https://nexss.info/buffet-6-in-september-24th-and-25th-2026-registration-opened/}\sep Online\\
2026 & \textbf{Co-chair and instructor}\sep Idealised modelling with LFRic: training course\sep Exeter, UK\\
2026 & \textbf{SOC}\sep UK Exoplanet Community Meeting~\link{https://ukexom2026.blogs.bristol.ac.uk/}\sep Bristol, UK\\
2025 & \textbf{SOC}\sep Atmospheric and interior evolution of planetary magma oceans~\link{https://www.lorentzcenter.nl/atmospheric-and-interior-evolution-of-planetary-magma-oceans.html}\sep Leiden, NL\\
2025 & \textbf{Co-chair}\sep BUFFET-5~\link{https://nexss.info/buffet-5-in-september-15th-and-16th-2025-registration-opened}\sep Bordeaux, FR\\
2025 & \textbf{SOC}\sep Exoclimes VII~\link{https://exoclimes.org/index_exovii.html}\sep Montreal, Canada\\
2025 & \textbf{Main organiser}\sep Idealised modelling with LFRic\sep Exeter, UK\\
2024 & \textbf{Co-chair}\sep BUFFET-4: Building a Unified Framework For Exoplanet Treatments~\link{https://nexss.info/buffet-4-in-october-22nd-and-23rd-2024/}\sep Online\\
2024 & What's Cookin' Doc? A CUISINES meeting (\ul{Chair})\sep Leiden, the Netherlands\\
2023 & \textbf{Co-chair}\sep ExoSLAM Summer School\sep Exeter, UK \\
2023 & \textbf{LOC}\sep Exoclimes VI\sep Exeter, UK \\
2023 & \textbf{Co-organiser}\sep Challenge of Science Leadership Short Course\sep Exeter, UK\\
2022 & Exoclimatology Theory Group Summer Retreat\sep Mawgan Porth, UK\\
\end{tabularx}


\section{Reviewing and Academic Service}
\begin{tabularx}{\linewidth}{@{}r X@{}}
Journals & Nat. Astron., MNRAS, Planet. Sci. J., Geophys. Res. Lett., ApJ, Planet. Space Sci., Q. J. R. Meteorol. Soc.\\
Funding  & STFC Consolidated Grant, STFC ERF, NASA XRP, STFC Small Award \\
Observations & James Webb Space Telescope General Observer Programs (Exoplanets \& Disks, Cycles 3 \& 4)\\
Membership & Royal Astronomical Society, Europlanet Society\\
\end{tabularx}


\section{Impact and Outreach}
\begin{tabularx}{\linewidth}{@{}r X@{}}
2026 & Public lecture \sep Bristol Astronomical Society \\
2026 & Astronomy on Tap \sep The Lansdown at Clifton, Bristol~\link{https://astronomyontap.org/event/astronomy-on-tap-bristol-exoplanets-weather-and-eclipses} \\
2023 & Interview by the University of Exeter about my research~\link{https://youtu.be/_Aus842EZWk}\\
2023 & Interview by UKRI/STFC~\link{https://youtu.be/PzQSfCwp0B8}\\
2023 & Expert Scientist at the British Science Festival Climate Exhibition~\link{https://exetersciencecentre.org/exhibition-2023}\\
2022 & Press releases: University of Exeter~\link{https://www.exeter.ac.uk/research/news/articles/questtouncoverintricacies.html}, American University~\link{https://www.american.edu/news/20220811-trappist-1.cfm}, \& INSU CNRS~\link{https://www.insu.cnrs.fr/fr/cnrsinfo/la-quete-pour-percer-le-mystere-des-climats-de-planetes-similaires-la-terre-avance}\\
2020--now & 3D visualisations of exoplanet simulations: \\
& `Cloudy Skies of Distant Exoplanets'~\link{https://issuu.com/universityofexeter/docs/images_of_research_2023} \sep University of Exeter Images of Research 2023\\
& `A refined look at tidally locked exoplanets'~\link{https://dirac.ac.uk/research-image-competition-2023} \sep DiRAC HPC Research Image Competition 2023\\
& `Exoplanetary Atmospheres'~\link{https://exetersciencecentre.org/gallery/denis_sergeev_uoe_exoplanets} \sep Exeter Science Centre, Science as Art Gallery 2020\\
& `Dusty exoplanet atmospheres'~\link{https://phys.org/news/2020-06-presence-airborne-signify-habitability-distant.html} \sep Nature Press Release\\
& `Virtual Reality Exploration of Exoplanets'~\link{https://youtu.be/Vnke9dWofR8} \sep 360 VR video (contributor)\\
2019 & Science consulting on the `Exoplanet Explorers' videogame\\
2015 & Blogging: \\
& Disastrous Disaster Movies~\link{http://www.scisnack.com/2015/12/17/disastrous-disaster-movies}\\
& Polar Lows: What Fuels Arctic Hurricanes?~\link{http://www.scisnack.com/2015/03/04/polar-lows-what-fuels-arctic-hurricanes}\\
& Worldwide Weird Weather Words~\link{http://www.scisnack.com/2014/12/17/worldwide-weird-weather-words}\\
\end{tabularx}


% \section{Technical Skills}
% \begin{tabularx}{\linewidth}{@{}l X@{}}
% % Languages & English (fluent), French (basic), Russian (native) \\
% Numerical models & LFRic \sep Unified Model \sep SOCRATES \sep LAGRANTO \\
% Programming languages & Python \sep FORTRAN \sep MATLAB \sep NCL \\
% Python libraries (user) & cartopy \sep cython \sep iris \sep matplotlib \sep numpy \sep pandas \sep pyvista \sep xarray \\
% Python libraries (creator/contributor) & aeolus \sep cartopy \sep pyvista \sep geovista \\
% Parallel computing & Dask \sep MPI \sep OpenMP \\
% Version control & Git \sep Subversion \\
% Document preparation & \LaTeX \sep Quarto \sep Jupyter Notebooks \sep Markdown \sep HTML \sep CSS \sep reST
% \end{tabularx}


\section{Vocational Training}
\begin{tabularx}{\linewidth}{@{}l X@{}}
2025--2026 & Postgraduate Certificate in Academic Practice (PGCAP) \sep University of Bristol \\
Sep 2023 & Belbin Training~\link{https://www.belbin.com} \\
Mar 2023 & Challenge of Science Leadership~\link{http://www.scienceleadership.co.uk} \\
Dec 2022 & Interview Training \\
Jul 2020 & Writing Workshop for Climate Scientists \\
Mar 2020 & ESA JWST Master Class~\link{https://ers-imaging.github.io/uk_workshop} \\
Jul 2019 & ICTP Summer School on Convective Organization and Climate Sensitivity~\link{https://indico.ictp.it/event/8669} \\
Apr 2018 & Fortran Modernisation Workshop~\link{https://www.nag.com/content/fortran-modernization-workshop} \\
Jan 2018 & Helicopter Underwater Escape Training Course (CA-EBS)~\link{http://www.petans.co.uk/courses/survival/huet-caebs/} \\
Dec 2017 & Sea Survival Course \\
Jun 2017 & Weather Presenting \\
Feb 2017 & Level 1 First Aid for Field Work Course \\
Jan 2017 & Raspberry Pi Basics \\
Apr 2016 & WWRP/WCRP/Bolin Center Polar Prediction School \\
Dec 2014 & UK Met Office Unified Model Training \\
% Dec 2011 & Global Climate Change course~\link{https://doi.org/10.1175/BAMS-D-11-00048.1} \\
\end{tabularx}


\section{Vocational Experience}
\begin{tabularx}{\linewidth}{@{}l X@{}}
Apr--Jun 2018 & Data Technician \\ & Processing of meteorological data collected in the IGP field campaign~\link{https://github.com/IGPResearch} \sep University of East Anglia \\
2015--2018 & Founder and Leader \\ & Python Users Group~\link{https://ueapy.github.io/} \sep University of East Anglia \\
Feb--Mar 2018 & Member of the Meteorology Team \\ & The Iceland-Greenland Seas Project (IGP) field campaign \sep Akureyri, Iceland \\
% Characterising the atmospheric forcing and the ocean response of coupled atmosphere-ocean processes; in particular cold-air outbreaks in the vicinity of the marginal ice zone and their effect on the ocean
Mar 2015 & Rapporteur \\ & Dynamics of Atmosphere-Ice-Ocean Interactions in the High-Latitudes~\link{https://highlatdynamics.w.uib.no} \sep Rosendal, Norway \\
Oct 2013 & Research Intern \\ & Geophysical Institute \sep University of Bergen, Norway \\
% & Supervisor: Alexey Eliseev \\
Aug--Sep 2013 & Trainee Forecaster \\ & Forecast and Briefing Service \sep Main Aviation Meteorological Centre, Vnukovo Airport \\
% Aug 2012 & Field practice in meteorology \sep Study of prevailing mesoscale processes via wind characteristic measurements and lake hydrothermodynamical modelling \sep Kamchatka, Russia \\
Jul 2012 & Research Intern \\ & A.M. Obukhov Institute of Atmospheric Physics \sep Moscow, Russia \\
% Jan--Feb 2012 & Field practice in meteorology \sep Karel Republic, Russia \\
\end{tabularx}


\section{Unsuccessful Funding Applications}
\begin{tabularx}{\linewidth}{@{}lX@{}}
2026 & GW4 Building Communities Generator Fund \\
2025 & STFC Small Award \\
2025 & Leverhulme Research Leadership Award \\
2025 & ERC Starting Grant \\
2025 & 2nd GPU Embedded CSE (eCSE) support call \\
2024 & 1st GPU Embedded CSE (eCSE) support call \\
2023 & STFC Ernest Rutherford Fellowship \\
2021 & Leverhulme Trust Early Career Fellowship \\
\end{tabularx}

\end{document}
